% Standards-mode starter template for ximeraExam 1.7.1.
% Keep ximeraExam.sty beside this file and compile with LuaLaTeX twice.
\DocumentMetadata{lang=en-US,pdfstandard=ua-2,tagging=on}
\documentclass{ximera}
\usepackage{unicode-math}
\usepackage[standards]{ximeraExam}

% Edit these details and replace the example standard codes with your own.
\examtitle{Standards Assessment}
\examcourse{Course Name}
\examdate{Exam Date}
\examversion{Form A}

% mark standard codes in the question text with \markStandard{<code>}.
% just the standard code itself will appear in the standard table, but the marked version will appear
% in the question text.  For example, \markStandard{ALG1} will appear as *ALG1* in the question text.
\RenewDocumentCommand{\markStandard}{m}{*#1*}

% Keep this for student copies; switch to \printanswers for the answer key.
\noprintanswers

\begin{document}
% build the cover page with a standard table instead of a grade table.  The standard table will show the standards
% that appear on each question.
\begin{examcover}
  \examcoverheading{\examtitletext}
  \noindent\examcoursetext\quad\examversiontext\quad\examdatetext\par
  \noindent\examnameprompt\ \examnamefield\par\bigskip
  \noindent\textbf{Instructions:} Show your work and justify your answers.
  \par\bigskip
  \centering
  \standardstable[v][noearned][4]
\end{examcover}

\begin{questions}

% A question may carry one or more standard codes.
\question[ALG1] Solve \(2x+3=11\). Show your work.
\begin{solution}[1.5in]
  Subtract \(3\), then divide by \(2\): \(x=4\).
\end{solution}

% An unmarked parent question inherits its parts' standards.
\question Consider the function \(f(x)=x^2\).
\begin{parts}
  \part[\markStandard{ALG2}] Evaluate \(f(3)\).
  \begin{solution}[0.75in]
    \(f(3)=9\).
  \end{solution}

  \part[ALG3] Solve \(f(x)=9\).
  \begin{solution}[\stretch{1}]
    \(x=-3\) or \(x=3\).
  \end{solution}
\end{parts}

\end{questions}
\end{document}
